\section{The entire tensor partner and a maximal-natural-degree tail-two gate}
\label{sec:entire-partner-support}

\paragraph{Status and exact scope.}
HD-FT18--19, reserved question Q2dr, is a derived consequence of the
accepted local support bounds and the actual tail-two gate.  The
tensor-normalizer argument was supplied independently in working
messages by \texttt{/root/imprimitive}; the full synthesis and the
maximal-natural-degree consequence are written here by
\texttt{/root/g34\_eightfolds}.  Complete author and separate root whole-proof readings pass, with the source/provenance/scope record read in full. No independent origin of the shared tensor-normalizer ingredient is claimed.  No new finite
certificate is needed or claimed for this universal argument.
No complete classification of possible
tensor partners is asserted.

For an operator \(S\) on a space \(E\), of finite order up to scalar, write
\[
 \nu(S)=\min_{z\in\mathbf C^*}\operatorname{rank}(S-zI)
       =\dim E-\max_z\dim\ker(S-zI).
 \tag{FTP.1}
\]
The quantity is unaffected by rescaling the operator.  Every operator
in a finite projective group can be rescaled to finite order, so is
semisimple and has the finite spectral distribution used below.

\subsection*{HD-FT18. A support floor independent of the number of partner factors}

Suppose the \emph{whole} partner space has a specified factorization
\[
 U=\bigotimes_{j=1}^s U_j,\qquad d_j=\dim U_j\ge2,
 \qquad b=\dim U=\prod_{j=1}^s d_j,
 \tag{FTP.2}
\]
and its actual finite projective action permutes the embedded full
matrix algebras \(\operatorname{End}(U_j)\otimes I\).
In particular there is no additional unbounded multiplicity factor
outside (FTP.2).  Fix positive numbers \(\eta_j\) such that every
compatible nonscalar operator on the \(j\)-th factor, whenever the
global operator fixes all factor algebras, has support at least
\(\eta_jd_j\).  Define
\[
                 \eta=\min\{1/4,\eta_1,\ldots,\eta_s\}.
 \tag{FTP.3}
\]

\begin{lemma}[HD-FT18]
Every nonscalar operator \(S\) in this actual projective
tensor-normalizer action satisfies \(\nu(S)\ge\eta b\).
Consequently, if all \(d_j\le A\), with \(A\ge2\), one may use
\[
                         \eta=1/\max\{A,4\}.
 \tag{FTP.4}
\]
The number of factors is arbitrary.  A genuine nonidentity scalar
on the whole source has actual moving rank equal to its full degree;
it is not assigned the projective support in (FTP.1).
\end{lemma}

\begin{proof}
An automorphism of a full complex matrix algebra is conjugation by
an invertible matrix.  Thus, after removing the actual permutation
of the factor algebras, an implementing operator is a tensor product
of factor operators, up to a scalar.  This is the matrix-algebra
normalizer argument already proved in D-BI19 and used in D-ATT02.
Factors in a permutation cycle have equal dimension.

If the operator fixes every factor algebra and is nonscalar, at
least one factor operator is nonscalar.  The maximum atom of a
convolution of finite probability distributions is at most the
maximum atom of each individual distribution.  Applying this to
the normalized spectral multiplicities gives
\(\nu(S)/b\ge\eta_j\) at that active factor.

Suppose instead there is a cycle of length \(\ell\ge2\), on
factors of degree \(d\ge2\).  Conjugate by factorwise operators so
that the cycle has only its cycle product as a twist.  This product
has finite projective order: a scalar power of the global operator,
after a multiple of the cycle length, makes the corresponding power
of every factor product scalar.  This follows because a scalar tensor
operator acts trivially by conjugation on every factor matrix algebra.
The cycle product can therefore be diagonalized after a scalar
normalization.  In a tensor-word basis, the resulting operator is
a weighted cyclic permutation.  On each word orbit an eigenvalue
has multiplicity at most one, since the successive nonzero weights
determine an eigenvector from its first coordinate.

There are exactly \(d\) constant words.  Every other word orbit
has length at least two, so every eigenspace on this cycle has
dimension at most
\[
 d+\frac{d^\ell-d}{2}
       =\frac{d^\ell+d}{2}\le\frac34d^\ell.
 \tag{FTP.5}
\]
The last inequality uses \(d^{\ell-1}\ge2\).  Convolution with
the remaining tensor cycles cannot increase the largest normalized
eigenspace.  Hence \(\nu(S)/b\ge1/4\).

Finally every nonscalar semisimple operator in degree \(d_j\)
has an eigenspace of codimension at least one, so
\(\eta_j=1/d_j\ge1/A\) is always available in the bounded-degree
case.  This proves (FTP.4) and the lemma.  For a genuine scalar,
the asserted full actual moving rank is immediate.
\end{proof}

Both local constants have actual operator witnesses.  A factorwise
diagonal involution with one negative line has normalized support
\(1/d\); a swap on \(\mathbf C^2\otimes\mathbf C^2\) has
eigenspace dimensions three and one, hence support \(1/4\).
These are operator statements, not asserted sphere-cover sources.

\subsection*{HD-FT19. Consequences for the actual tail-two source}

Retain exactly the actual source of HD-FT13--15: a connected sphere
cover, an actual quotient onto \(P=A_n\) or \(S_n\), and an actual
stable factorization
\[
 V=W\otimes U\quad\hbox{up to scalars},\qquad
 W=S^{(n-2,2)}\ \hbox{or}\ S^{(n-2,1,1)},\qquad
 e=\dim V,\qquad r\ge12,\qquad
 \sum_{i=1}^r\operatorname{rank}(1-T_i)=2e+6.
 \tag{FTP.6}
\]
The transposed labels and irreducible ordinary alternating
restrictions are permitted.  Every original inertia is genuine;
the projected tuple is not assigned twelve genuine marks.
The whole partner has the actual tensor-normalizer structure (FTP.2).

For \(n\ge10^6\), accepted HD-GC07 and HD-FT15 give at least four
genuine entries with identity projection to \(P\), and their total
normalized actual moving rank is at most \(61/n^2\).  If instead
one uses the qualitative Monderer--Neftin pair classification at
\(n>N_{14}\), with \(n\ge64\), the same conclusion has at least
seven such entries.  Put \(q=4\) in the effective case and \(q=7\)
in that separately scoped qualitative case.

\begin{theorem}[HD-FT19]
Under these actual hypotheses,
\[
                             q\eta\le\frac{61}{n^2}.
 \tag{FTP.7}
\]
In particular, if the entire partner consists of factors of degree
at most \(A\), independently of their number, then
\[
                     qn^2\le61\max\{A,4\}.
 \tag{FTP.8}
\]

More generally, suppose every partner factor is of one of the
following two kinds:
\begin{itemize}
\item a factor of degree at most \(A\), where \(A\ge2\);
\item a standard or ordinary tail-two representation of \(A_m\),
where \(15\le m\le n\), whose actual simple-core action is
normalized by every compatible local factor operator.
\end{itemize}
The actual image must still permute the full factor matrix algebras
and preserve the indicated simple-core actions when it fixes those
factors.  Then
\[
                  qn^2\le61\max\{4,A,n-1\}.
 \tag{FTP.9}
\]
For the effective case \(q=4\), \(n\ge10^6\), this forces
\[
                             A\ge\frac{4n^2}{61}.
 \tag{FTP.10}
\]
In particular there is no source in which the distinguished actual
stable tail-two factor has maximal natural degree and the entire
partner consists only of standard and tail-two factors, regardless
of the number of those factors.  Bounded additional factors are
also excluded once \(n^2>61A/4\).

Every conclusion also permits additional partner factors, of any
degree, for which the entire compatible local normalizer has a
proved normalized projective support floor of at least \(1/4\).
This is an explicit support hypothesis, not an assertion about an
unclassified residual factor.
\end{theorem}

\begin{proof}
At an identity-projection entry, \(T_i\) is a scalar multiple of
\(I_W\otimes S_i\).  If \(S_i\) is nonscalar, its actual rank is
at least \(e\nu(S_i)/b\ge\eta e\), with every scalar phase
allowed.  If \(S_i\) is scalar, genuineness of \(T_i\) instead
gives full rank \(e\ge\eta e\).  Summing over the at least \(q\)
entries and using the accepted upper budget proves (FTP.7).
Equation (FTP.8) follows from (FTP.4).

For the ordinary factors, \(\operatorname{Aut}(A_m)=S_m\) at
\(m\ge15\).  Faithfulness, irreducibility, and the chosen
ordinary extension show that every compatible local normalizer
is a scalar times the corresponding \(S_m\) operator, as in
D-ATT02 and D-AST02.  On the standard module, projective
faithfulness and integral codimension give a floor \(1/(m-1)\).
The exact D-ATT01 minima give respectively
\[
 \frac{\nu(S^{(m-2,2)}(x))}{m(m-3)/2}\ge\frac2m,
 \qquad
 \frac{\nu(S^{(m-2,1,1)}(x))}{(m-1)(m-2)/2}
                                  \ge\frac2{m-1}
 \tag{FTP.11}
\]
for every nonidentity \(x\in S_m\).  Both exceed or equal
\(1/(m-1)\ge1/(n-1)\).  The identity permutation is a scalar
local operator and is already handled separately.  Sign choices
of extension do not affect projective support.

Thus HD-FT18 gives
\(\eta\ge1/\max\{4,A,n-1\}\), proving (FTP.9).
Additional factors with the stated quarter-support floor do not
decrease this value, by the definition (FTP.3).
For \(n\ge10^6\), one has \(4n^2>61(n-1)\) and
\(4n^2>244\), so neither \(n-1\) nor four can satisfy the
maximum in (FTP.9) with \(q=4\).  This forces (FTP.10).
When no bounded-degree pieces are present, one may put \(A=2\);
the resulting contradiction proves the stated maximal-natural-degree
case.  If \(\dim U=1\), there are no nonscalar partner factors,
but all \(q\) entries have full rank, immediately contradicting
\(q\le61/n^2\).  This also covers the empty tensor product.
\end{proof}

\paragraph{Strictness and retained boundary.}
The earlier partner-dimension gate alone says
\(\dim U\ge4n^2/61\) in the effective case.  It does not prevent
\(\dim U=2^s\) with arbitrarily many two-dimensional factors.
Equations (FTP.7)--(FTP.8) exclude that entire actual tensor-normalizer
class at \(n\ge10^6\), since its floor is \(1/4\), independently
of \(s\).  The displayed dimensions can satisfy the old numerical
inequality, but no cover or branch tuple is inferred from them.
This is a strict narrowing of that previously unrestricted
large-partner family under explicit decomposition hypotheses.

The theorem does not assert that every possible partner admits such
a complete factorization, that all remaining factors have bounded
degree, or that an unbounded residual multiplicity space can be
discarded.  It also does not choose a distinguished stable factor
from a group that permutes it with the rest: the actual quotient
and factorization (FTP.6) remain indispensable.  The exact
Monderer--Neftin threshold \(N_{14}\) remains unevaluated; the
effective and qualitative counts retain their separate ranges.
